%% get rid of autoindent   :setl noai nocin nosi inde=
%%
%% 1. pdflatex FSI
%% 2. bibtex FSI
%% 3. pdflatex FSI
%% 3. pdflatex FSI
%\documentclass[review]{elsarticle}
\documentclass[]{article}
%\documentclass[preprint]{elsarticle}

\usepackage{lineno,hyperref}
\modulolinenumbers[5]

\usepackage[fleqn]{amsmath} 
\usepackage{amssymb}
\usepackage{graphicx}
\usepackage{tabularx} 
\usepackage{footnote}
\usepackage{titlesec}
\titleformat{\section}{\large\bfseries}{\thesection}{1em}{}

\makesavenoteenv{tabular}
\makesavenoteenv{table}
\usepackage[percent]{overpic}
\usepackage{subfigure}% subcaptions for subfigures
\usepackage{subfigmat}% matrices of similar subfigures, aka small multiples
\usepackage{array}
\setlength{\mathindent}{0pt}

\renewcommand{\labelitemi}{$\triangleright$}
\renewcommand{\labelitemii}{$\triangleright$}
\renewcommand{\labelitemiii}{$\triangleright$}
\renewcommand{\labelitemiv}{$\triangleright$}


%\journal{Journal of Computational Physics}

\newcommand{\mb}{\mathbf}
\newcommand{\tn}{\textnormal}
\newcommand{\unit}[1]{\ensuremath{\, \mathrm{#1}}}

\newcommand{\bmz}{\mbox{\boldmath $z$\unboldmath}}
\newcommand{\bmr}{\mbox{\boldmath $r$\unboldmath}}
\newcommand{\bmb}{\mbox{\boldmath $b$\unboldmath}}
\newcommand{\bms}{\mbox{\boldmath $s$\unboldmath}}
\newcommand{\bmg}{\mbox{\boldmath $g$\unboldmath}}
\newcommand{\bmq}{\mbox{\boldmath $q$\unboldmath}}
\newcommand{\bmU}{\mbox{\boldmath $U$\unboldmath}}
\newcommand{\bmu}{\mbox{\boldmath $u$\unboldmath}}
\newcommand{\bmv}{\mbox{\boldmath $v$\unboldmath}}
\newcommand{\bmw}{\mbox{\boldmath $w$\unboldmath}}
\newcommand{\bmm}{\mbox{\boldmath $m$\unboldmath}}
\newcommand{\bmi}{\mbox{\boldmath $i$\unboldmath}}
\newcommand{\bmj}{\mbox{\boldmath $j$\unboldmath}}
\newcommand{\bmk}{\mbox{\boldmath $k$\unboldmath}}
\newcommand{\bmx}{\mbox{\boldmath $x$\unboldmath}}
\newcommand{\bmX}{\mbox{\boldmath $X$\unboldmath}}
\newcommand{\bmH}{\mbox{\boldmath $H$\unboldmath}}
\newcommand{\bmy}{\mbox{\boldmath $y$\unboldmath}}
\newcommand{\bmn}{\mbox{\boldmath $n$\unboldmath}}
\newcommand{\bmf}{\mbox{\boldmath $f$\unboldmath}}
\newcommand{\bmF}{\mbox{\boldmath $F$\unboldmath}}
\newcommand{\bmS}{\mbox{\boldmath $S$\unboldmath}}
\newcommand{\bmG}{\mbox{\boldmath $G$\unboldmath}}
\newcommand{\bmV}{\mbox{\boldmath $V$\unboldmath}}
\newcommand{\bmI}{\mbox{\boldmath $I$\unboldmath}}
\newcommand{\bmD}{\mbox{\boldmath $D$\unboldmath}}
\newcommand{\bmnu}{\mbox{\boldmath $\nu$\unboldmath}}

\newcommand{\DefTen}{\mathbb{D}}
\newcommand{\EyeTen}{\mathbb{I}}

\DeclareMathOperator{\Ei}{Ei}

\newcommand{\MVBcmnt}[1]{{\color{red}{\em #1}}}

\newcolumntype{L}[1]{>{\raggedright\let\newline\\\arraybackslash\hspace{0pt}}m{#1}}
\newcolumntype{C}[1]{>{\centering\let\newline\\\arraybackslash\hspace{0pt}}m{#1}}
\newcolumntype{R}[1]{>{\raggedleft\let\newline\\\arraybackslash\hspace{0pt}}m{#1}}

%%\bibliographystyle{elsarticle-num}
\bibliographystyle{acm}
%%%%%%%%%%%%%%%%%%%%%%%

\title{Hybrid AI CFD for simulating atomization and spray in
 internal combustion engines}

%%!!format authors properly
%\author{
%	LastName1, FirstName1\\
%	\texttt{first1.last1@xxxxx.com}
%	\and
%	LastName2, FirstName2\\
%	\texttt{first2.last2@xxxxx.com}
%}
\author{
  Sussman, Mark \\
  Dept. of math, FL State Univ. 
  \and
  Shoele, Kourosh  \\
  Dept. of Mech. Engr., FL State Univ.
  \and
  Arienti, Marco \\
  Sandia National Laboratory, Livermore CA.
}

\begin{document}
\maketitle

ABSTRACT:
A new AI accelerated algorithm is proposed for simulating atomization and
spray in internal combustion engines which are problems of national 
interest.  Applications are the improvement in rotating detonation
engines and military jet engines.
 
 \linenumbers

\section{Mathematical Modeling}
There are $M$ materials and each material $1\le m \le M$ occupies the region,
\begin{eqnarray}
	\Omega_{m}(t)=\left\{ (\vec{x},t) | \phi_{m}(\vec{x},t)\ge 0 \right\}
\end{eqnarray}
The $\phi_{m}(\bmx,t)$ are level set 
functions\cite{sussman1994level,osher1988fronts}.  
The zero
level set of such functions demarcate the boundary of material 
$m$:
\begin{eqnarray}
	\partial\Omega_{m}(t)=\left\{ (\vec{x},t)) | \phi_{m}(\vec{x},t)=0 
	\right\}
\end{eqnarray}
In the bulk regions of material $m$ we have the following
equations for compressible materials:
\begin{eqnarray}
	\rho_{t}+\nabla\cdot(\rho\vec{u})=0
\end{eqnarray}
\begin{eqnarray}
	(\rho\vec{u})_{t}+\nabla\cdot(
	\rho\vec{u}\otimes\vec{u}+p\EyeTen-\tau)=\vec{F}_{body}
\end{eqnarray}
\begin{eqnarray}
	(\rho E)_{t}+\nabla\cdot(
	(\rho E + p)\vec{u}-\vec{u}\cdot\tau-k\nabla T)=
	\vec{u}\cdot \vec{F}_{body}
	\label{energy}
\end{eqnarray}
\begin{eqnarray}
	(\rho Y)_{t}+\nabla\cdot(\rho Y\vec{u})=\nabla\cdot(\rho D\nabla Y)
\end{eqnarray}

\begin{eqnarray}
	\tau=2\mu\bar{D}
\end{eqnarray}
\begin{eqnarray}
	\bar{D}=D-\frac{\mbox{Trace}(D)\EyeTen}{\mbox{Trace}(\EyeTen)}
\end{eqnarray}
\begin{eqnarray}
	D=\frac{\nabla\vec{u}+(\nabla\vec{u})^{T}}{2}
\end{eqnarray}
\begin{eqnarray}
	\rho E=\rho e_{\mbox{int}}+
	\frac{1}{2}\rho|\vec{u}|^{2}
\end{eqnarray}
\begin{eqnarray}
	p=p(\rho,e) \hspace{10pt} e=c_{v}T
\end{eqnarray}
If the material is a viscoelastic FENE-CR material, then one has the
additional equation for the configuration tensor $A$,
%StewartLaySussmanOhta
\begin{eqnarray}
	A_{t}+\vec{u}\cdot\nabla A=A\cdot(\nabla\vec{u})^{T}+
	\nabla\vec{u}\cdot A-\frac{f(R)}{\lambda}(A-\EyeTen),
\end{eqnarray}
and the momentum constitutive law term is modified,
\begin{eqnarray}
	\tau=\tau+\frac{\tilde{\eta}_{P}}{\lambda}f(R)A.
\end{eqnarray}
If the material is an elastic material undergoing plastic deformation, then
one has the following additional equation for the deviatoric stress
tensor,
\begin{eqnarray}
	s_{t}+\vec{u}\cdot\nabla s+s\Omega-\Omega s=2G(\bar{D}-D^{p}),
\end{eqnarray}
and the momentum constitutive law term is modified,
\begin{eqnarray}
	\tau=\tau+s.
\end{eqnarray}
If the material is an elastic material undergoing plastic deformation, then
the plastic strain equation is solved too,
\begin{eqnarray}
	\epsilon^{p}_{t}+\vec{u}\cdot\nabla\epsilon^{p}=
	\sqrt{\frac{2}{3}}|\Lambda|.
\end{eqnarray}
\begin{eqnarray}
	|\Lambda|=\sqrt{\frac{3}{2}}\sqrt{D^{p}:D^{p}}
\end{eqnarray}
\begin{eqnarray}
	D^{p}=\Lambda N  \hspace{10pt} N=\sqrt{3/2}\frac{s}{S_{e}}
	\hspace{10pt} (\dot{\epsilon}^{p})^{2}=
	\frac{2}{3}\Lambda^{2}
	\label{plastic_strain_rate}
\end{eqnarray}
$\Omega$ is the spin tensor,
\begin{eqnarray}
	\Omega_{ij}=\frac{1}{2}(
	\frac{\partial u_{i}}{\partial x_{j}}-
	\frac{\partial u_{j}}{\partial x_{i}})
\end{eqnarray}
The effective stress is defined as,
\begin{eqnarray}
	S_{e}=\sqrt{\frac{3}{2}}\sqrt{s:s}
\end{eqnarray}
The Johnson-Cook model yield stress is given by,
\begin{eqnarray}
	\sigma_{y}=
	(A+B(\epsilon^{p})^{n})
	(1+C\ln(
	\frac{\dot{\epsilon}^{p}}
	{\dot{\epsilon}_{0}}
	))(1-\theta^{m}) \hspace{10pt}
	\theta=\frac{T-T_{0}}{T_{m}-T_{0}}
	\label{Johnson}
\end{eqnarray}
One form of the Zerilli Armstrong model yield stress is given by,
\begin{eqnarray}
	\beta=\beta_{0}-\beta_{1}
	\log(\dot{\epsilon}^{p}/\dot{\epsilon}^{p}_{ref})
\end{eqnarray}
\begin{eqnarray}
	\sigma_{y}=A+B(\epsilon^{p})^{n}e^{-\beta T}
	\label{zerilli}
\end{eqnarray}

The energy contribution due to elasto-plastic deformation is modeled as,
\begin{eqnarray}
	\beta\dot{\epsilon}^{p}S_{e},
\end{eqnarray}
and is added to (\ref{energy}).  In order to determine $\Lambda$, as it
appears in (\ref{plastic_strain_rate}), the following ``radial
return algorithm'' is implemented 
(see \cite{tran2004particle} for the case when $n=1$ in (\ref{Johnson})
or (\ref{zerilli})):
\begin{description}
\item[1.] $\dot{\epsilon}^{p}=0$, $k=0$
\item[2.] $S_{e}=\sqrt{\frac{3}{2}}\sqrt{s:s}$
\item[3.] $\sigma_{y}=\sigma_{y}(T,\epsilon^{p},\dot{\epsilon}^{p})$
\item[4.]
	if $S_{e}-\sigma_{y}\le 0$ then $\Lambda=0$ (i.e. $D^{p}:D^{p}=0$).
\item[5.]
	if $S_{e}>\sigma_{y}$ then,
\item[6.]
	Solve the following equation for $\epsilon^{p}_{new}$:
        \begin{eqnarray}
		S_{e}-\sigma_{y}=3G(\epsilon^{p}_{new}-\epsilon^{p})+
		B((\epsilon^{p}_{new})^{n}-
		  (\epsilon^{p})^{n})
        \end{eqnarray}
\item[7.] 
	 \begin{eqnarray}
		 \dot{\epsilon}^{p}=\frac{
			 \epsilon^{p}_{new}-\epsilon^{p}}{\Delta t}
	 \end{eqnarray}
\item[8.] 
	 \begin{eqnarray}
		 \Lambda=\sqrt{\frac{3}{2}}\dot{\epsilon}^{p}
	 \end{eqnarray}
	 \begin{eqnarray}
		 D^{p}=\Lambda N
	 \end{eqnarray}
\item[9.] $k=k+1$
\item[10.] if $k=1$ then go back to step 3.
\end{description}

The mathematical model at sharp interfaces is given as follows.  First,
define a ``closest point'' function:
\begin{eqnarray}
	\vec{x}_{c}(\vec{x})=\mbox{argmin}_{\vec{y},m}
	\left\{ ||\vec{y}-\vec{x}|| | \phi_{m}(\vec{y})=0 \right\}
\end{eqnarray}
Taking into account mass transfer, a source material's level set function
(e.g. ice that is melting or liquid that is turning into vapor) satisfies
\begin{eqnarray}
	(\phi_{m_{s}})_{t}+\vec{u}_{m_{s}}(\vec{x}_{c}(\vec{x}))\cdot
	\nabla\phi_{m_{s}}=
	-\frac{|\dot{m}(\vec{x}_{c}(\vec{x}))|}{\rho_{m_{s}}}
	|\nabla\phi_{m_{s}}|
\end{eqnarray} 
The destination material's level set function satisfies
\begin{eqnarray}
	(\phi_{m_{d}})_{t}+\vec{u}_{m_{s}}(\vec{x}_{c}(\vec{x}))\cdot
	\nabla\phi_{m_{d}}=
	\frac{|\dot{m}(\vec{x}_{c}(\vec{x}))|}{\rho_{m_{s}}}
	|\nabla\phi_{m_{d}}|
\end{eqnarray} 
The above equations are equivalent to the following:
\begin{eqnarray}
	(\phi_{m_{s}})_{t}+\vec{u}_{m_{d}}(\vec{x}_{c}(\vec{x}))\cdot
	\nabla\phi_{m_{s}}=
	-\frac{|\dot{m}(\vec{x}_{c}(\vec{x}))|}{\rho_{m_{d}}}
	|\nabla\phi_{m_{s}}|
\end{eqnarray} 
\begin{eqnarray}
	(\phi_{m_{d}})_{t}+\vec{u}_{m_{d}}(\vec{x}_{c}(\vec{x}))\cdot
	\nabla\phi_{m_{d}}=
	\frac{|\dot{m}(\vec{x}_{c}(\vec{x}))|}{\rho_{m_{d}}}
	|\nabla\phi_{m_{d}}|
\end{eqnarray} 
Now, we define the interface normal:
\begin{eqnarray}
	\vec{n}_{m}=\frac{\nabla\phi_{m}}{||\nabla\phi_{m}||}
\end{eqnarray}
The interface curvature is,
\begin{eqnarray}
	\kappa_{m}=\nabla\cdot\vec{n}_{m}
\end{eqnarray}
The rate of mass transfer is,
\begin{eqnarray}
	\dot{m}=\frac{
		k_{m_{s}}\nabla T_{m_{s}}-
		k_{m_{d}}\nabla T_{m_{d}}}{L}\cdot \vec{n}_{m_{d}}.
	\label{heat_flux}
\end{eqnarray}
The phase change interface temperature condition is,
\begin{eqnarray}
	T_{m_{s}}=T_{m_{d}}=
	T_{\mbox{saturation}}(P_{\mbox{equilibrium}})
\end{eqnarray}
For evaporation, the following rate of mass transfer should match that in
(\ref{heat_flux}):
\begin{eqnarray}
	\dot{m}=\frac{-\vec{n}_{m_{s}}\cdot\nabla Y_{m_{d}}\rho_{m_{d}}D}
	{Y^{\Gamma}-1}
\end{eqnarray}
%FIX ME!
Let the tensor $\sigma$ represent pressure, viscosity, viscoelastic, and
elastic contributions to the momentum equations.  One has:
\begin{eqnarray}
	\sigma=-p\EyeTen+\tau.
\end{eqnarray}

The stress jump conditions are
\begin{eqnarray}
	\vec{n}_{m_{s}}\cdot
	(-\sigma_{m_{s}}+\sigma_{m_{d}})
	\cdot\vec{n}_{m_{s}}=
	-\gamma_{m_{s},m_{d}}\kappa_{m_{s}}-
	\dot{m}^{2}(
	\frac{1}{\rho_{m_{s}}}-
	\frac{1}{\rho_{m_{d}}})
\end{eqnarray}
\begin{eqnarray}
	(\vec{t}_{i,tangent})_{m_{s}}\cdot
	(-\sigma_{m_{s}}+\sigma_{m_{d}})
	\cdot\vec{n}_{m_{s}}=0 \hspace{10pt} i=1,2
\end{eqnarray}
The velocity jump conditions are
\begin{eqnarray}
	(
	\vec{u}_{m_{s}}-
	\vec{u}_{m_{d}})\cdot\vec{n}_{m_{s}}=
	\dot{m}(
	\frac{1}{\rho_{m_{d}}}-
	\frac{1}{\rho_{m_{s}}})
\end{eqnarray}

\section{budget}

The budget should be arount 495K for this phase I grant.

\begin{description}
	\item[1.] Sussman; two months (~28K)
	\item[2.] Shoele; ?
	\item[3.] Arienti (100K)
        \item[4.] graduate students
	\item[5.] postdoc
\end{description}

\section{Applications}

\begin{description}
	\item[1.] Jet in cross flow problem (military jet engines, Pratt
		and Whitney)
	\item[2.] Diesel injectors
	\item[3.] Rotating Detonation Engines (see latest article with Everett
		Wenzel and Marco Arienti).
\end{description}

\section{Proposed algorithm}

\begin{description}
	\item[0.] Reseed particles with dense particle regions in narrow
		bands about interfaces (but not on interfaces!).  
		Initialize the state of the particles
		using the associated PINN from the previous time step.
	\item[1.] Cell integrated semi Lagrangian advection and
		particle advection. Particle advection uses PINN velocity.
	\item[2.] Cell integrated semi Lagrangian advection 
		due to phase 
		change.  
	\item[3.] Viscosity, Viscoelastic, and elastic forces.
	\item[4.] Temperature Diffusion
	\item[5.] Mass Fraction Diffusion
	\item[6.] Pressure Gradient
	\item[7.] Identify all separate contiguous material regions (this is
		done already, see Wang, Simakhina, Sussman, or Sussman Ohta)
	\item[8.] For each separate region, train a NN using PINN, CFD, and
		experimental data.
		The CFD data on the interface of each material region has the
		most weight.  The PINN is trained in a Lagrangian frame of 
		reference.  $x,t^{n+1}$ mapped to itself in the NN. 
		$x-u\Delta t,t^{n}$ mapped to $x,t^{n}$ in the NN.  The NN is
		valid for $t^{n}<t<t^{n+1}$.  It could be that trained NN
		can be recycled if the associated material regions 
		are topologically similar.
	\item[9.] Go back to step 0.
\end{description}

%%\biboptions{sort&compress}
\bibliography{cmof_paper}

\end{document}
%%% Local Variables:
%%% mode: latex
%%% TeX-master: t
%%% End:
